\documentclass[10pt]{article}
\usepackage[margin=12mm]{geometry}
\usepackage{graphicx}
\usepackage{caption}
\pagestyle{empty}
\captionsetup{font=small,labelfont=bf}

\begin{document}
\begin{figure}[p]
\centering
\begin{minipage}[t]{0.485\textwidth}
  \centering
  \includegraphics[width=\linewidth]{fig5A_ce_canonical_definition.pdf}
\end{minipage}\hfill
\begin{minipage}[t]{0.485\textwidth}
  \centering
  \includegraphics[width=\linewidth]{fig5B_ce_canonical_major_subtrees.pdf}
\end{minipage}

\vspace{2mm}
\includegraphics[width=\textwidth]{fig5C_ce_canonical_all_subtrees.pdf}

\caption{\textbf{Canonical positions and lineage/null proximity on subtree Pareto fronts.}
%
\textbf{(A)} Null-independent canonical coordinates. For assignment $a$,
travel distance is normalized as
$D_1(a)=[C_T(a)-C_T(T)]/[C_T(S)-C_T(T)]$ and cell-state distance as
$D_2(a)=[C_S(a)-C_S(S)]/[C_S(T)-C_S(S)]$, the resulting coordinates place
the travel and cell-state optima at $T=(0,1)$ and $S=(1,0)$, respectively.
$P^*$ is the sampled Pareto assignment closest to the natural lineage $L$ in
this geometry. The horizontal coordinate $u$ is the arc length from $T$ to
$P^*$ divided by the total Pareto front arc length from $T$ to $S$, and
$d_{LP}=\left\|D(L)-D(P^*)\right\|_2$ is the endpoint-normalized
distance between the natural lineage and the sampled Pareto assignment. For
the first-cousin-null mean $N$, $d_{NP}=\left\|D(N)-D(P^*)\right\|_2$
uses that same lineage-selected $P^*$. Distances are evaluated only at actual
sampled assignments, not interpolated line segments. Neutral gray denotes
$u$, purple denotes $d_{LP}$, and gold denotes the first-cousin reference and
$d_{NP}$.
%
\textbf{(B)} Lineage-to-front and first-cousin-null-to-front distances for
the major subtrees P0, AB, ABa, ABp, and P1. Each subtree retains the same
marker shape at its purple $d_{LP}$ and gold $d_{NP}$ positions; direct labels
identify the subtree.
%
\textbf{(C)} Canonical positions of all 43 subtrees containing at least 12 usable terminal cells.
Horizontal position records the travel and cell-state compromise of $P^*$;
vertical position is the natural lineage's endpoint-normalized distance
$d_{LP}$ to that assignment. Gold markers add $d_{NP}$ references for the five
major subtrees only. Marker shape identifies the root or first-order lineage
region, with AB separated from the root marker. ABplpp(p) denotes the nearly
overlapping, directly nested ABplpp and ABplppp records. The endpoint-normalized
coordinates of $L$ used to calculate $d_{LP}$ are retained without clipping to
$[0,1]$.}
\label{fig:ce_canonical_summary}
\end{figure}
\end{document}
